% Canonical node 008 · C004A; current section path 1.2.1.
% Canonical owner: C004A
\cnnaNodeHeading{008}{C004A}{First provenance slot $s_1$}
\cnnaNodePlacement{1.2.1}{D1}

\paragraph*{Position and scientific role.}
\texttt{C004A} is the first node in the bootstrap group.  It consumes the
finite $b$-ary provenance grammar from \texttt{C003} and the canonical
breadth-first/lexicographic order from \texttt{C018}.  Its sole task is to
select the first \emph{structural provenance slot}.  It does not perform a
birth, add a vertex to the carrier, assign a time, or introduce conductance,
response or geometry.

% CNNA-ARCHITECTURE-BEGIN C004A
\begin{cnnaInfoBox}{CNNA handoff: selecting the first open complement slot}
C004A selects the least admissible continuation of the realized root without yet turning it into a node.  It is therefore the first explicit handoff from an open complement possibility to a birth input.  C013, not C004A, owns the actual realization and relation.
\end{cnnaInfoBox}
% CNNA-ARCHITECTURE-END C004A

\paragraph*{Explicit construction.}
Let $\varepsilon$ denote the C003 root address and let the active branching
domain satisfy $b\ge2$.  The least local C003 rank is the zero-based slot
$q_0=0\in S_b$.  The first provenance slot is therefore
\begin{equation}
  s_1:=\bigl(\varepsilon,q_0,\varepsilon\mathbin{\Vert}q_0\bigr),
  \qquad
  \operatorname{parent}(s_1)=\varepsilon,
  \quad
  \operatorname{rank}(s_1)=0,
  \quad
  \operatorname{addr}(s_1)=(0).
  \label{eq:c004a-first-slot-construction}
\end{equation}
The subscript in $s_1$ is a one-based ordinal name: ``the first slot.''  It
does not change the stored zero-based sibling rank.

The formal object carries exactly the two direct predecessors, a C003 grammar
$G$ and a C018 schedule $\Sigma$, together with the coherence condition
$\Sigma.\mathrm{grammar}=G$.  Parent, rank and address are derived functions,
not independently supplied fields.

\paragraph*{Structural invariants and minimality.}
The address has depth one and C003 recovers both of its syntactic components:
\begin{equation}
  |\operatorname{addr}(s_1)|=1,
  \qquad
  \operatorname{parent?}(\operatorname{addr}(s_1))=\varepsilon,
  \qquad
  \operatorname{finalSlot?}(\operatorname{addr}(s_1))=0.
  \label{eq:c004a-recovery}
\end{equation}
Rank zero is the least local slot.  C018 therefore places $s_1$ before every
distinct sibling of the root.  More strongly, for every non-root provenance
address $c$,
\begin{equation}
  c=(0)
  \quad\lor\quad
  (0)\prec_{\mathrm{BFSlex}}c.
  \label{eq:c004a-global-minimum}
\end{equation}
For depth-one addresses this is the increasing-rank branch of C018; for depth
at least two it follows from breadth-first depth priority.  Thus $(0)$ is the
unique minimum non-root address under the already-proved strict order.  This
is a theorem about the structural address space, not yet about the first
accepted birth event.

\paragraph*{Cutoff boundary.}
Finite admission is deliberately separated from structural existence.  Define
\begin{equation}
  \operatorname{WithinCutoff}(s_1)
  \quad\Longleftrightarrow\quad
  |\operatorname{addr}(s_1)|\le L.
\end{equation}
Since the depth is exactly one,
\begin{equation}
  L\ge1\Longrightarrow\operatorname{WithinCutoff}(s_1),
  \qquad
  L=0\Longrightarrow\neg\operatorname{WithinCutoff}(s_1).
  \label{eq:c004a-cutoff}
\end{equation}
At $L=0$ the intrinsic word $(0)$ and the label $s_1$ remain well-defined,
but the finite approximant is root-only and C018 enumerates no admissible
birth slot.  This prevents the cutoff from being mistaken for part of the
local branching grammar.

\paragraph*{Cross-layer realization.}
Python stores the two predecessors in an immutable dataclass and computes
\texttt{parent}, \texttt{rank}, \texttt{address} and cutoff admission as
properties.  It rejects predecessor objects of the wrong type and rejects a
C003 grammar paired with a C018 schedule built from another grammar.  Lean
stores the same predecessors in \texttt{FirstProvenanceSlot} and internalizes
their coherence as the proof field \texttt{schedule\_grammar}.  The Lean
lemmas \texttt{address\_parent} and \texttt{address\_finalSlot} prove the
syntactic recovery equations; \texttt{address\_eq\_or\_before\_nonroot} and
\texttt{address\_before\_distinct\_nonroot} prove global minimality; the two
cutoff theorems prove the exact $L=0$ versus $L\ge1$ boundary.

\paragraph*{Counterchecks and ownership boundary.}
The Python suite checks the one-based/zero-based distinction, all active sample
values $b\in\{2,3,5\}$ and $L\in\{0,1,3\}$, first-position agreement with the
C018 enumeration for positive cutoff, rejection of incoherent predecessors,
and absence of geometry, node IDs, event indices, times, conductance and
response fields.  These tests exclude both an off-by-one slot and accidental
promotion of a provenance label into a physical or dynamical state.

\paragraph*{Result and downstream handoff.}
\texttt{C004A} closes the canonical structural selection
\begin{equation}
  s_1=(\varepsilon,0,(0)),
\end{equation}
with $(0)$ minimal among non-root addresses and admitted exactly when $L\ge1$.
Edges \texttt{E005} and \texttt{E064} are the verified incoming handoffs from
\texttt{C003} and \texttt{C018}.  Verified edge \texttt{E007} supplies this
slot to \texttt{C013}, where the first non-root vertex and first directed
relation are constructed.  No such birth is owned by the present node.

\noindent\textbf{Primary code anchors.}
Python:
\path{derivation/code/python/src/cnna/derivation/s01_primitive_response_coupled_finite_provenance/s02_bootstrap_of_the_first_relation/s01_c004a__first_provenance_slot_s1.py},
class and derived properties, lines 36--82; constructors, lines 85--99; tests,
lines 19--100.  Lean:
\path{derivation/code/lean/core/src/CNNA/Derivation/S01_PrimitiveResponseCoupledFiniteProvenance/S02_BootstrapOfTheFirstRelation/S01_C004A_FirstProvenanceSlotS1.lean},
predecessor structure and construction, lines 23--45; address invariants and
minimality, lines 46--168; cutoff boundary, lines 169--185.  The Supplement
lists all 23 declaration-level anchors and their exact source hashes.
